% ============================================================
\chapter{Future Scope}
\label{sec:future}
% ============================================================

The most pressing work is \textbf{closing the security gaps}: binding trusted
signing keys to the identities that authenticated a link, requiring TLS rather
than preferring it, checking signatures on every import path, and authorising
users more finely than by one shared secret. Next to it, \textbf{a relay of last
resort} would join the peer pairs that no direct path can reach, without making
it the default data path.

Placement would gain from \textbf{locality and topology}: accounting for which
node already holds the input data, as iDDS \cite{idds2025} does, and for
network distance through a graph-structured view of the cluster
\cite{zhao2023}. \textbf{Fault tolerance for the submitter and the lead} would
let a parallel script outlive its submitter and a training ring outlive rank~0's
machine, and \textbf{a cluster-wide view} would merge every machine's events in
one place.

The system can reach more hardware and more programs. \textbf{More accelerators}
means AMD GPUs through ROCm and phones that run whole jobs; \textbf{memory
escalation} on out-of-memory would try gradient accumulation, checkpointing and
lower precision before failing; and \textbf{additional runtimes} would extend
interception beyond Python.

For training, \textbf{fast synchronisation for larger rings} could build on
compressed exact synchronisation with PowerSGD, which kept accuracy at 4 to 11
ranks in this work, and on synchronising the model in staggered fragments
(Streaming DiLoCo \cite{douillard2025streaming}). Finally, \textbf{a controlled
multi-machine evaluation} across separate machines and networks, including a
machine with two GPUs, would measure speedup where the container lab cannot.
